\documentclass[dvipdfmx,twoside]{jsarticle}
\usepackage{amsmath,amssymb}
\usepackage{CJKutf8}
\usepackage{okumacro}
\usepackage{booktabs}
\usepackage{array}
\usepackage{xcolor}
\usepackage{colortbl}
\usepackage{tipa}
\usepackage{tikz}
\usepackage{fancybox}
\usepackage{wasysym}
\usepackage{pifont}
\usepackage{graphicx}
\usepackage{wrapfig}
\usepackage{geometry}
\geometry{margin=2cm}
\usepackage{fancyhdr}
\usepackage[utf8]{inputenc}
\usepackage{CJK}
% 设置页面样式
\pagestyle{fancy}
\fancyhf{}
\renewcommand{\headrulewidth}{0pt}
\fancyhead[RO]{数学–\thepage}
\fancyhead[LE]{数学–\thepage}

% 重新定义 plain 样式
\fancypagestyle{plain}{%
  \fancyhf{}%
  \renewcommand{\headrulewidth}{0pt}%
  \fancyhead[RO]{数学-\thepage}%
  \fancyhead[LE]{数学-\thepage}%
}
% 自定义A方框命令
\newcommand{\abb}[1]{%
\begin{tikzpicture}[baseline]
\node[draw=black, 
      rectangle, 
      minimum width=0.8cm, 
      minimum height=0.3cm, 
      fill=gray!25, 
      font=\bfseries,
      line width=1pt,
      inner sep=2pt,
      anchor=base] {#1};
\end{tikzpicture}%
}
\newcommand{\ab}[1]{%
\begin{tikzpicture}[baseline]
\node[draw=black, 
      rectangle, 
      minimum width=0.8cm, 
      minimum height=0.3cm, 
      font=\bfseries,
      line width=1pt,
      inner sep=2pt,
      anchor=base] {#1};
\end{tikzpicture}%
}
\newcommand{\sab}[1]{%
\begin{tikzpicture}[baseline]
\node[draw=black, 
      rectangle, 
      minimum width=0.4cm, 
      minimum height=0.05cm, 
      font=\bfseries,
      line width=1pt,
      inner sep=0.5pt,
      anchor=base] {#1};
\end{tikzpicture}%
}

\newcommand{\maru}[1]{\tikz[baseline=-0.7ex]{
    \node[shape=circle,draw,inner sep=1pt,minimum size=5pt,anchor=center] {\footnotesize #1};}}
\definecolor{headercolor}{RGB}{220,220,220}
\definecolor{rowcolor1}{RGB}{245,245,245}
\definecolor{rowcolor2}{RGB}{255,255,255}

% \title{\vspace{-1.5cm} 2025年8月羚課文科数学月考 }
% \author{\textnosfal{Linc\ -\ 伊}}
\date{}
\begin{document}

\begin{CJK}{UTF8}{ipxm}  % 使用ipxm字体
  %月考表纸部分
\begin{center}

\vspace*{5cm}

% 学校logo
\includegraphics[width=5cm]{pics/1.jpg}

\vspace{2cm}

% 主标题
{\fontsize{24}{30}\selectfont\bfseries\sffamily
2025年10月\\
\vspace{1em}
羚課文科数学月考
}

\end{center}
\newpage
%第一部分
\noindent
\begin{tikzpicture}
\node[draw, rectangle, minimum width=1cm, minimum height=0.8cm, font=\Huge] {I};
\end{tikzpicture}
\\
\\
\textbf{問1}\qquad $P = 12a^2 + 14ab - 21bc - 18ca$ とする。

\vspace{0.8em}

(1) \quad $P$ を因数分解すると

\vspace{0.5em}

\[
P = \left( \ab{\color{red}{\textsf{A}}} a + \ab{\color{red}{\textsf{B}}} b \right) \left( \ab{\color{red}{\textsf{C}}} a - \ab{\color{red}{\textsf{D}}} c \right)
\]

\vspace{0.3em}

である。

\vspace{1.2em}

(2) \quad $6a = \sqrt{6}$, $14b = \sqrt{2} + \sqrt{3} - \sqrt{6}$, $18c = \sqrt{12} - \sqrt{8}$ とすると

\vspace{0.5em}

\[
P = \frac{\ab{\color{red}{\textsf{E}}} + \ab{\color{red}{\textsf{F}}} \sqrt{\ab{\color{red}{\textsf{G}}}}}{\ab{\color{red}{\textsf{H}}}}
\]
\newpage
\begin{center}
- 計算欄 (memo) -
\end{center}
\newpage

\noindent
\textbf{問2}\qquad 2次関数 \( f(x) = -x^2 + 4x + 5 \) を考える。
\begin{enumerate}
\item[(1)]
放物線 \( y = f(x) \) の頂点の座標は
\( \left(\ab{\color{red}{\textsf{I}}}\, ,\,\ab{\color{red}{\textsf{J}}} \right) \)
である。\\[2em]
\item[(2)]
放物線 \( y = f(x) \) を \( x \) 軸方向に \( k \), \( y \) 軸方向に \(-4\) だけ平行移動して得られる放物線を
\( y = g(x) \) とすると
\[
g(x) = -\left( x -\ab{\color{red}{\textsf{K}}} - k \right)^2 +\ab{\color{red}{\textsf{L}}}
\]
である。
\end{enumerate}

\newpage
\begin{center}
- 計算欄 (memo) -
\end{center}
\newpage
%第二部分
\noindent
\begin{tikzpicture}
\node[draw, rectangle, minimum width=1cm, minimum height=0.8cm, font=\Huge] {II};
\end{tikzpicture}
\\
\\
\noindent
2つの袋A、Bがある。Aの袋には白カードが4個，赤カードが1個入っており，Bの袋には白カードが2個，赤カードが3個入っている。はじめにAの袋から同時に2個のカードを取り出し，続いて，Bの袋から同時に2個のカードを取り出す。

\vspace{1.5em}

(1) \quad Aから2個の白カードを取り出し，Bからは白カードと赤カードをそれぞれ1個ずつ取り出す確率は$\dfrac{\ab{\color{red}{\textsf{A}}}}{\ab{\color{red}{\textsf{BC}}}}$ \quad 

\vspace{1.5em}

(2) \quad 取り出した4個のカードの中に，3個の白カードと1個の赤カードが入っている確率は \quad $\dfrac{\ab{\color{red}{\textsf{D}}}}{\ab{\color{red}{\textsf{E}}}}$ \quad である。

\vspace{1.5em}

(3) \quad 取り出した4個のカードがすべて同じ色である確率は \quad $\dfrac{\ab{\color{red}{\textsf{F}}}}{\ab{\color{red}{\textsf{GH}}}}$ \quad である。

\vspace{1.5em}

(4) \quad 取り出した4個のカードの中に含まれる白カードが2個以下である確率は \quad $\dfrac{\ab{\color{red}{\textsf{IJ}}}}{\ab{\color{red}{\textsf{KL}}}}$ \quad である。\\

% \textbf{問2}\qquad 2つの2次関数

% \[
%   f(x) = \frac{1}{3}x^2 + ax - b, \qquad g(x) = -x^2 + cx + b
% \]

% が、次の2つの条件 (A), (B) を満たすような $a$, $b$, $c$ を求めよう。

% \begin{itemize}
%   \item[(A)] グラフ $y=f(x)$ とグラフ $y=g(x)$ は2つの直線 $x=-1,~x=3$ 上で交わっている。\\[0.5em]
%   \item[(B)] $g(x)$ の最大値と $f(x)$ の最小値の差は $\dfrac{16}{3}$ である。
% \end{itemize}
% \ \\[0.5em]
% 条件(A) より2つの2次関数のグラフは直線 $x=-1$ 上で交わっているから
% \[
%   3a + \ab{\color{red}{\textsf{M}}} b - \ab{\color{red}{\textsf{N}}} c = \ab{\color{red}{\textsf{O}}} \qquad \cdots\ (1)
% \]
% であり、また、直線 $x = 3$ 上で交わっているから
% \[
%   \ab{\color{red}{\textsf{P}}} a - 2b - \ab{\color{red}{\textsf{Q}}} c = -12 \qquad \cdots\ (2)
% \]
% を得る。①と②より $b = \ab{\color{red}{\textsf{R}}}$ である。\\[0.5em]

% \noindent
% 次に、条件(B) と $b = \abb{\color{red}{\textsf{R}}}$ より\\[0.5em]
% \[
%   \ab{\color{red}{\textsf{S}}} a^2 + 3c^2 = \ab{\color{red}{\textsf{TU}}}
% \]
% を得る。これらから\\[0.5em]
% \[
%   a = -\frac{\ab{\color{red}{\textsf{V}}}}{\ab{\color{red}{\textsf{W}}}}, \qquad
%   c = \ab{\color{red}{\textsf{X}}}
% \]
% となる。

\newpage
\begin{center}
- 計算欄 (memo) -
\end{center}
\newpage
%第三部分
\noindent
\begin{tikzpicture}
\node[draw, rectangle, minimum width=1cm, minimum height=0.8cm, font=\Huge] {III};
\end{tikzpicture}
\\

2つの整数 \(a = 588, b = 1260\) を考える。

\begin{enumerate}
\item
\(a, b\) を素因数分解すると
\[
a = 2^{\sab{\textsf{\color{red}{A}}}} \cdot 3 \cdot {\ab{\textsf{\color{red}{B}}}}^{\sab{\textsf{\color{red}{C}}}}, \quad b = 2^{\ab{\textsf{\color{red}{D}}}} \cdot 3^{\ab{\textsf{\color{red}{E}}}} \cdot \ab{\textsf{\color{red}{F}}} \cdot \ab{\textsf{\color{red}{G}}}
\]
である。ただし, \(\abb{\textsf{\color{red}{F}}} < \abb{\textsf{\color{red}{G}}}\) とする。よって, \(a, b\) の最大公約数は \(\ab{\textsf{\color{red}{HI}}}\) である。\\[1em]

\item
次の条件 (i), (ii) を満たす正の整数 \(c\) を考える。
\begin{itemize}
  
    \item[(i)] \(a, b, c\) の最大公約数は, \(a, b\) の最大公約数に等しい。
    \item[(ii)] \(a, b, c\) の最小公倍数は, \(a, b\) の最小公倍数の2倍である。
\end{itemize}
条件(i), (ii)を満たすような\(c\)は全部で \(\ab{\textsf{\color{red}{J}}}\) 個ある。そのような中で最小のものは \(\ab{\textsf{\color{red}{KLM}}}\) である。\\[1.3em]

\item
整数 \(a, b\) に対して, 方程式
\[
ax - by = \abb{\textsf{\color{red}{KLM}}} \quad \cdots (1)
\]
の整数解 \(x, y\) を求めよう。
方程式 (1) の整数解 \(x, y\) を求めるためには, 方程式
\[
\ab{\textsf{\color{red}{N}}}x - \ab{\textsf{\color{red}{OP}}}y = \ab{\textsf{\color{red}{Q}}} \quad \cdots (2)
\]
の整数解 \(x, y\) を求めればよい。
方程式 (2) を満たすような正の整数 \(x, y\) で, \(y\) が最も小さいものを求めると
\[
x = \ab{\textsf{\color{red}{RS}}}, \quad y = \ab{\textsf{\color{red}{T}}}
\]
である。したがって, 方程式 (1) の整数解は
\[
x = \ab{\textsf{\color{red}{UV}}} + \ab{\textsf{\color{red}{WX}}}k, \quad y = \ab{\textsf{\color{red}{Y}}} + \ab{\textsf{\color{red}{Z}}}k 
\]
である。ただし, \(k\) は整数である。
\end{enumerate}

\newpage
\begin{center}
- 計算欄 (memo) -
\end{center}
\newpage
%第四部分
\noindent
\begin{tikzpicture}
\node[draw, rectangle, minimum width=1cm, minimum height=0.8cm, font=\Huge] {IV};
\end{tikzpicture}
\\

下の右図は，四面体 $OABC$ の展開図である．四面体 $OABC$ において
\[
BC=10, \quad AC=8, \quad \sin\angle ACB=\frac{3}{4}
\]
\[
OA=3, \quad \triangle ABC \equiv \triangle OBC
\]
が成り立つとする．

\vspace{1.5em}

\begin{minipage}{0.48\linewidth}
\centering
\begin{tikzpicture}[scale=2]

% ---------- 左图 ----------
% 定义点
\coordinate (A) at (0,0);
\coordinate (B) at (0.4,-0.8);
\coordinate (C) at (1.8,0);
\coordinate (O) at (0.6,1);

% 绘制多边形
\draw (A)--(B)--(C)--(O)--cycle;
\draw[solid] (O)--(B);
% 绘制虚线
\draw[dashed] (A)--(C);

% 标记点
\node[left]  at (A) {$A$};
\node[below] at (B) {$B$};
\node[right] at (C) {$C$};
\node[above] at (O) {$O$};

% ---------- 右图 ----------
\begin{scope}[xshift=4.5cm] % 平移以分开两图

% 定义点
\coordinate (B) at (0,0);
\coordinate (C) at (2,0);
\coordinate (O1) at (0.2,1.35);
\coordinate (O2) at (1.5,1.5);
\coordinate (O3) at (0.85,-1);
\coordinate (A) at (0.85,1);

% 绘制多边形
\draw (B)--(C)--(A)--cycle;
\draw (B)--(O3)--(C);
\draw (B)--(O1);
\draw (C)--(O2);
\draw (A)--(O1);
\draw (A)--(O2);

% 绘制对角线
\draw (A)--(B);
\draw (A)--(C);

% 标记点
\node[above] at (O1) {$O$};
\node[above] at (O2) {$O$};
\node[below] at (O3) {$O$};
\node[left]  at (B) {$B$};
\node[right] at (C) {$C$};
\node[above] at (A) {$A$};

\end{scope}
\end{tikzpicture}
\end{minipage}

\vspace{2em}

\begin{enumerate}
\item 三角形 $ABC$ の面積は　$\ab{\color{red}{\textsf{AB}}}$　である。

\vspace{1em}
\item 点 $A$ から辺 $BC$ におろした垂線を $AH$ とすると，$AH$ の長さは
　$\ab{\color{red}{\textsf{C}}}$　である。

\vspace{1em}
\item 平面 $ABC$ と平面 $OBC$ のなす角を $\theta$ とすると
\[
\cos\theta = \cfrac{\ab{\color{red}{\textsf{D}}}}{\ab{\color{red}{\textsf{E}}}},\qquad
\sin\theta = \dfrac{\sqrt{\ab{\color{red}{\textsf{FG}}}}}{\ab{\color{red}{\textsf{H}}}}
\]
である。

\vspace{1em}
\item 四面体 $OABC$ の体積は　
$\cfrac{\ab{\color{red}{\textsf{IJ}}} \sqrt{\ab{\color{red}{\textsf{KL}}}}}{\ab{\color{red}{\textsf{M}}}}$　である。
\end{enumerate}
\newpage
\begin{center}
- 計算欄 (memo) -
\end{center}

\end{CJK}

\end{document}